\section{Theoretische Grundlagen}
\label{sec:grundlagen}

Der Aufbau eines \emph{Digitalen Zwillings} setzt Kenntnisse aus Fabrikplanung, Virtueller Inbetriebnahme und industrieller Kommunikationstechnik voraus.
Die im Folgenden erarbeiteten Grundlagen bilden den begrifflichen Rahmen, auf den Entwicklung, Inbetriebnahme und Bewertung der in dieser Arbeit realisierten \emph{Digitalen Zwillinge} in den nachfolgenden Kapiteln aufbauen.

Das vorliegende Kapitel erarbeitet das methodische und technische Fundament dafür.
Die klassische Fabrikplanung nach VDI~5200 strukturiert den Planungsprozess in
sequenzielle Phasen, deren Grenzen durch proprietäre Datenformate und fehlende
Interoperabilität der eingesetzten Werkzeuge entstehen.
Das Konzept der Digitalen Fabrik nach VDI~4499 ergänzt dieses Modell durch ein
durchgängiges Datenmanagement, das diese Phasengrenzen überwindet und die bisher
isoliert arbeitenden Planungsabschnitte miteinander vernetzt.
Anschließend werden die Virtuelle Inbetriebnahme nach VDI/VDE~3693 sowie
Klassifikationen Digitaler Zwillinge eingeführt.
Das Referenzarchitekturmodell RAMI~4.0 (DIN~SPEC~91345) und der normative Rahmen
nach ISO~23247 verorten den \emph{Digitalen Zwilling} im industriellen Kontext,
bevor abschließend die eingesetzten Technologien und Kommunikationsstandards
vorgestellt werden.

\subsection{Fabrikplanung und Digitale Fabrik}
Die Richtlinie VDI~5200 Blatt~1 gliedert die Fabrikplanung in sieben sequenzielle Phasen,
von der \emph{Zielfestlegung} (Phase~1) über \emph{Grundlagenermittlung},
\emph{Konzeptplanung} und \emph{Detailplanung} bis hin zu
\emph{Realisierungsvorbereitung}, \emph{Realisierungsüberwachung} und
\emph{Hochlaufbetreuung} (Phase~7).
Damit liefert sie das methodische Grundgerüst für die systematische Planung von
Produktionsanlagen \cite{vdi5200}.
Sämtliche Phasen werden von einem übergreifenden Projektmanagementprozess begleitet,
wie Abbildung~\ref{fig:vdi5200_phasen} verdeutlicht.

\begin{figure}[H]
    \centering
    \includegraphics[trim={30pt 655pt 30pt 40pt},clip,width=\linewidth]{vdi5200_phasenmodell}
    \caption{Phasenmodell der Fabrikplanung nach VDI~5200 Blatt~1, Bild~1
        (Quelle: \cite{vdi5200}).}
    \label{fig:vdi5200_phasen}
\end{figure}

Das Phasenmodell ist dabei grundsätzlich sequenziell ausgerichtet, wobei spätere Phasen
auf den Ergebnissen früherer aufbauen. Klassischerweise werden die einzelnen Phasen
jedoch mit weitgehend eigenständigen, kaum miteinander verknüpften Modellen und
Werkzeugen bearbeitet.
Um diese Grenzen zwischen den Phasen zu überwinden, setzt die Digitale Fabrik auf ein
durchgängiges Datenmanagement, das alle Planungsphasen miteinander vernetzt.

\begin{figure}[H]
    \centering
    \includegraphics[width=\linewidth]{vdi4499_phasenmodell}
    \caption{Digitaler Fabrikbetrieb in den Lebenszyklusphasen eines Produktionssystems
        (Quelle: VDI~4499 Blatt~2 \cite{vdi4499b2}, Bild~2).}
    \label{fig:vdi4499_phasen}
\end{figure}

Das gezeigte Bild verdeutlicht, wie VDI~4499 Blatt~2 dieses Konzept in den
Lebenszyklusphasen eines Produktionssystems konkretisiert \cite{vdi4499b2}.
VDI~4499 Blatt~1 definiert die Digitale Fabrik als vernetztes System digitaler Modelle,
Methoden und Werkzeuge zur ganzheitlichen Planung und Steuerung von Fabrikprozessen
\cite{vdi4499} und schafft damit eine gemeinsame Datenbasis über alle Phasengrenzen hinweg.

\subsection{Virtuelle Inbetriebnahme nach VDI/VDE~3693}
\label{subsec:defizite}
Die \emph{Virtuelle Inbetriebnahme} (VIBN) setzt bereits in den Planungs- und
Realisierungsphasen an und ermöglicht es, Steuerungsprogramme und Produktionsabläufe
an einem digitalen Anlagenmodell zu erproben und zu optimieren, noch bevor die reale
Anlage fertiggestellt und die eigentliche Inbetriebnahme durchgeführt wird.
Die Richtlinie VDI/VDE~3693 Blatt~1 definiert VIBN als methodisches Vorgehen zur
frühzeitigen Absicherung der Steuerungssoftware vor der realen Inbetriebnahme einer
Anlage \cite{vdi3693} und unterscheidet dabei verschiedene Modellarten
(\emph{Software-in-the-Loop}, \emph{Hardware-in-the-Loop}).
Technische Grundlage bilden dabei cyber-physische Produktionssysteme (CPPS), die durch
bidirektionale Echtzeit-Kopplung zwischen physischer Anlage und digitalem Abbild die
Voraussetzung für Digitale Zwillinge schaffen \cite{monostori2014}.
Lechler et al.\ belegen empirisch, dass VIBN die Qualität und Effizienz im Engineering
steigert, Anlaufzeiten durch frühzeitige Fehlererkennung verkürzt und eine sicherere
Inbetriebnahme ermöglicht, indem erkannte Fehler bereits vor der realen Inbetriebnahme
behoben werden \cite{lechler2019}.

\subsection{Klassifikation Digitaler Zwillinge}
Um Digitale Zwillinge im Kontext dieser Arbeit präzise einordnen zu können, wird die
in der Literatur am weitesten verbreitete Klassifikation nach Kritzinger et al.
herangezogen.
Sie unterscheidet anhand des Datenflusses zwischen physischem Objekt und digitalem Abbild
drei Reifegrade \cite{kritzinger2018}.
\begin{itemize}
    \item \textbf{Digitales Modell:} Keine automatisierte Datenkopplung. Jegliche
        Aktualisierung zwischen physischem und digitalem Objekt erfolgt manuell.
    \item \textbf{Digitaler Schatten:} Ein automatisierter, unidirektionaler Datenfluss
        vom physischen zum digitalen Objekt. Das digitale Objekt spiegelt den Zustand
        des physischen wider, wirkt aber nicht auf dieses zurück.
    \item \textbf{Digitaler Zwilling:} Ein vollständig integrierter, in beide
        Richtungen automatisierter Datenfluss zwischen dem physischen und dem digitalen
        Objekt \cite{kritzinger2018}. Veränderungen am physischen Objekt führen
        unmittelbar zu Anpassungen des digitalen Abbildes und umgekehrt.
\end{itemize}
Diese Klassifikation bildet den Bewertungsmaßstab für den in den nachfolgenden Kapiteln
beschriebenen Umsetzungsgrad.

\iffalse
\subsubsection{Definition nach Stark und Damerau}
Stark und Damerau halten in der CIRP Encyclopedia of Production Engineering fest,
dass trotz intensiver Forschung bis heute keine einheitlich akzeptierte
wissenschaftliche Definition des Begriffs Digitaler Zwilling existiert.
Ihre Arbeit dokumentiert allein 19 verschiedene Definitionen aus der Literatur
\cite{stark2019}.
Den Ursprung des Konzepts führen sie auf Grieves zurück, der es 2002 an der University
of Michigan als \enquote{Mirrored Spaces Model} eingeführt hat.
Ihre eigene Definition lautet: Ein Digitaler Zwilling ist \enquote{ein digitales Abbild
eines aktiven, einzigartigen Produkts oder Produktionssystems, das durch ausgewählte
Eigenschaften, Bedingungen und Verhaltensweisen mittels Modellen, Informationen und
Daten charakterisiert wird} \cite{stark2019}.

Um zu verstehen, was das konkret bedeutet, unterscheiden Stark und Damerau drei
Bestandteile, die erst zusammen einen echten Digitalen Zwilling ergeben:

\textbf{(1) Digital Master:} Der \emph{Digital Master} ist das in der Entwicklungsphase
erstellte virtuelle Modell einer Anlage, das Geometrie, Kinematik und
Materialeigenschaften beschreibt.
Für den Maze Runner ist dies das aus CAD-Daten exportierte 3D-Modell in Isaac Sim,
das alle Bauteile originalgetreu abbildet.
Für den Dobot Magician ist es die kinematische Beschreibung des Roboterarms mit seinen
Gelenken, Drehachsen und Gelenklimits.

\textbf{(2) Digital Shadow:} Der \emph{Digital Shadow} sind die Echtzeitdaten, die die
physische Anlage laufend erzeugt: Sensorwerte, Schaltstellungen und Positionen.
Die Steuerungseinheit der physischen Anlage übermittelt diese Daten kontinuierlich an
das digitale Modell, sodass der aktuelle Betriebszustand stets im digitalen Abbild
sichtbar ist.

\textbf{(3) Verknüpfung:} Erst die Verbindung zwischen Konstruktionsmodell und
Livedaten erzeugt den eigentlichen Digitalen Zwilling.
Eine Middleware-Schicht liest die Betriebsdaten der physischen Anlage aus und setzt sie
in Bewegungsänderungen des digitalen Modells um, sodass beide Instanzen synchron laufen.

Besonders betonen Stark und Damerau den Aspekt der \emph{Einzigartigkeit}: Ein
Digitaler Zwilling bildet stets eine bestimmte physische Instanz ab, nicht ein
generisches Produktionslinienmodell.
\fi

\subsection{Digitale Zwillinge im Fertigungskontext}
Die Klassifikation nach Kritzinger et al.\ beschreibt, in welchem Maß der Datenfluss
zwischen physischem Objekt und digitalem Abbild automatisiert ist.
Sie beantwortet damit die Frage nach dem Kopplungsgrad, beschreibt aber nicht, welche
physischen Elemente beobachtet werden können und wie das Gesamtsystem dabei strukturiert ist.
Die Normenreihe ISO~23247 schließt diese Lücke und beschreibt das Digital-Twin-Framework
für die Fertigung als schichtbasiertes IoT-Rahmenwerk, das Abbildung~\ref{fig:dt_manufacturing}
veranschaulicht.

\begin{figure}[H]
    \centering
    \includegraphics[width=\linewidth]{DT for manufacturing}
    \caption{IoT-Rahmenwerk für Digitale Zwillinge in der Fertigung mit Anwendungsschicht,
        Digitalem Zwilling, Gerätekommunikation und beobachtbaren Fertigungselementen
        (Quelle: ISO/FDIS~23247-1 \cite{iso23247}, Figure~2).}
    \label{fig:dt_manufacturing}
\end{figure}

Das Framework gliedert sich in drei Schichten.
Die Anwendungsschicht umfasst Funktionen wie Echtzeitsteuerung, Offline-Analytik und
vorausschauende Wartung, die auf den Digitalen Zwilling zugreifen.
Der Digitale Zwilling bildet die mittlere Schicht und repräsentiert je nach Anwendungsfall
eines oder mehrere der acht beobachtbaren Fertigungselemente.
Diese umfassen Personal, Betriebsmittel, Material, Fertigungsprozesse,
Fertigungsinfrastruktur, Umgebung, Produkt und Begleitdokumente \cite{iso23247}.
Die Gerätekommunikation verbindet den Digitalen Zwilling bidirektional mit den physischen
Fertigungselementen und übernimmt den Datentransport zwischen realer Anlage und
digitalem Abbild.
Rechts der Abbildung sind zudem Nutzenpotenziale aufgeführt, die Digitale Zwillinge
in der Fertigung erschließen, darunter Echtzeitregelung, Produktionsplanung und
vorausschauende Risikobewertung \cite{iso23247}.
Die Darstellung vermittelt damit einen anschaulichen ersten Überblick über das
Zusammenspiel von digitalem Abbild, Gerätekopplung und Anwendungsnutzen.

\subsection{Referenzarchitekturmodell RAMI~4.0}
An Planung, Umsetzung und Betrieb vernetzter Industrie-4.0-Anlagen sind Fachleute
aus verschiedenen Disziplinen beteiligt. Jede Disziplin beschreibt Komponenten in
eigenen Begriffen, sodass Zuständigkeiten über Systemgrenzen hinweg schwer
abzustimmen und Probleme an Schnittstellen schwer einzugrenzen sind.
Für die Einordnung nach Betriebsebene diente traditionell die Automatisierungspyramide,
die Feldgeräte, Steuerungen und IT-Systeme in eine vertikale Hierarchie gliedert.
Sie erfasst jedoch weder den Produktlebenszyklus noch die Architekturschichten, die
für eine vollständige Beschreibung vernetzter Industrie-4.0-Komponenten notwendig sind.
Das Referenzarchitekturmodell RAMI~4.0 (DIN~SPEC~91345) schließt diese Lücke und
erweitert die Automatisierungspyramide um die Achsen Lebenszyklus und Wertschöpfung
sowie sechs Architekturschichten. Es stellt damit den dreidimensionalen Rahmen bereit,
in dem sich Maschinen, Anlagen, Softwarekomponenten und Digitale Zwillinge einheitlich
beschreiben und zueinander in Beziehung setzen lassen \cite{dinspec91345}.
Das Modell richtet sich an alle Akteursgruppen, die an Industrie-4.0-Systemen
mitwirken.

\begin{itemize}
    \item Maschinenbauer und Anlagenbauer (mechanische Konstruktion, Steuerungstechnik)
    \item Elektro- und Automatisierungsingenieure (SPS, Feldbusse, Sensorik)
    \item IT-Architekten und Softwareentwickler (Cloud, MES, ERP-Anbindung)
    \item Inbetriebnehmer vor Ort
    \item Betreiber und Produktionsplaner
    \item Systemintegratoren
    \item Normungsgremien und Zertifizierungsstellen
\end{itemize}

In Abbildung~\ref{fig:rami40} ist das Modell als dreidimensionale Struktur dargestellt.

\begin{figure}[H]
    \centering
    \includegraphics[width=\linewidth]{rami.jpg}
    \caption{RAMI~4.0: Referenzarchitekturmodell Industrie~4.0 als dreidimensionales
        Schichtenmodell mit den Achsen \enquote{Life Cycle \& Value Stream},
        \enquote{Hierarchy Levels} und \enquote{Layers}
        (Quelle: DIN~SPEC~91345 \cite{dinspec91345}).}
    \label{fig:rami40}
\end{figure}

Grundlegende Einheit des Modells ist das \textbf{Asset}, worunter jedes physische
oder virtuelle Objekt mit definierten Eigenschaften und Wert verstanden wird, vom
einzelnen Sensor über eine Steuerungseinheit bis zur gesamten Anlage.
Für jedes Asset kann eine Verwaltungsschale (Asset Administration Shell) angelegt
werden, die seine Eigenschaften und Fähigkeiten maschinenlesbar digital repräsentiert.

Die \textbf{Hierarchieachse} übernimmt die Struktur der Automatisierungspyramide mit
den Ebenen Feldgerät, Steuerungseinheit, Arbeitsstation und Arbeitszelle und ergänzt
sie um Betrieb sowie Vernetzte Welt, welche übergeordnete Cloud- und IoT-Dienste
einschließt. Jedes Asset lässt sich einer oder mehreren dieser Ebenen zuordnen.

In vertikaler Richtung gliedern sechs \textbf{Architekturschichten} jede Hierarchieebene
und beschreiben, welche Fähigkeiten ein Asset dort bereitstellt. Sie umfassen die Schichten Asset (die physische Komponente selbst),
Integration, Kommunikation, Information, Funktionen und Geschäftsprozesse.
Kommunikationsstandards wie OPC~UA sind der Kommunikationsschicht zuzuordnen,
Datenformate wie OpenUSD der Informationsschicht.

Die dritte Achse erfasst den \textbf{Lebenszyklus} eines Assets von der Entwicklung
über den Betrieb bis zur Außerbetriebnahme und unterscheidet dabei zwischen Typ, dem
allgemeinen Modell einer Komponente, und Instanz, dem konkreten Exemplar im Einsatz.
Digitale Zwillinge bilden stets Instanzen ab, da sie den aktuellen Zustand einer
bestimmten physischen Anlage repräsentieren.

Damit steht Maschinenbauern, Elektroingenieuren, IT-Architekten und Inbetriebnehmern
eine gemeinsame Beschreibungssprache zur Verfügung. Systeme lassen sich
schichtenweise analysieren, Schnittstellen eindeutig verorten und Probleme an
Systemgrenzen gezielt eingrenzen, ohne das Gesamtsystem vollständig durchleuchten
zu müssen.

\subsection{Normativer Rahmen nach ISO~23247}

Im vorangehenden Unterkapitel wurde das Digital-Twin-Framework in seiner Grundstruktur
dargestellt.
Eine formale Beschreibung der einzelnen Entitäten und ihrer Schnittstellen fehlt darin.
Für Fertigungsumgebungen liegt mit ISO~23247 ein internationaler Standard vor, der
diesen Rahmen beschreibt.
Sie ist in vier Teile gegliedert. Teil~1 beschreibt allgemeine Grundsätze und
Anforderungen, Teil~2 die Referenzarchitektur mit funktionalen Sichten, Teil~3
die Informationsattribute der beobachtbaren Fertigungselemente und Teil~4 die
technischen Anforderungen an den Informationsaustausch \cite{iso23247}.

Das in Teil~2 spezifizierte \emph{Digital Twin Framework} unterscheidet dabei zwei
funktional verschiedene Einsatzkontexte.
Im Unternehmenskontext sind Aspekte wie die Integration übergeordneter Systeme
(ERP, MES), rollenbasierte Zugangskontrolle sowie die Vernetzung mehrerer Zwillinge
untereinander über die Cross-System Entity bedeutsam.
Für einen Digitalen Zwilling als eigenständige Instanz, die eine einzelne physische
Anlage konsistent abbildet, stehen hingegen die Core Entity als Träger der digitalen
Repräsentation und die \emph{Data Collection and Device Control Entity} (DCDCE) als Kommunikationsschicht zur realen Anlage im Vordergrund.
Abbildung~\ref{fig:iso23247} zeigt das funktionale Referenzmodell des Frameworks mit
der hierarchischen Anordnung seiner vier Entitäten.

\begin{figure}[H]
    \centering
    \includegraphics[width=0.85\linewidth]{{Functional view of the Digital Twin reference model for manufacturing}}
    \caption{Funktionale Referenzarchitektur des Digital Twin Frameworks für die Fertigung
        mit User Entity, Core Entity, Data Collection and Device Control Entity
        und Cross-System Entity sowie den beobachtbaren Fertigungselementen außerhalb
        der Frameworkgrenze
        (Quelle: \cite{iso23247}).}
    \label{fig:iso23247}
\end{figure}

Ein Digitaler Zwilling, der eine einzelne Fertigungsanlage mit zugehöriger
Nutzerschnittstelle abbildet, ohne in Unternehmenssysteme eingebettet zu sein oder
auf weitere Zwillinge zuzugreifen, erfordert primär die direkte Kopplung zwischen
physischer und digitaler Welt.
Die \textbf{DCDCE} ist dabei die funktional zentrale Entität. Sie erfasst kontinuierlich
die Zustandsdaten der physischen Anlage, leitet sie an die Core Entity weiter und
übersetzt in Gegenrichtung Steuerbefehle in Signale für die physische Anlage.
Den Kommunikationsteil dieser Entität bildet eine Middleware-Schicht, die den
Datentransport zur physischen Anlage übernimmt. Protokolle wie OPC~UA oder ROS~2 werden
in dieser Rolle eingesetzt. Die DCDCE reicht jedoch über den reinen Datentransport
hinaus und umfasst auch die Steuerungslogik für Gerätebefehle sowie die
Synchronisationsschnittstelle zur Core Entity.

Die \textbf{Core Entity} strukturiert ihre Funktionen über mehrere Functional Interfaces.
Für einen eigenständigen Digitalen Zwilling sind die Schnittstellen für digitale
Modellbildung (\emph{Digital Modeling FI}), Datensynchronisation (\emph{Synchronization FI})
und physikalische Simulation (\emph{Simulation FI}) zentral.
Die \emph{Reporting FI} ermöglicht die strukturierte Ausgabe von Zustandsdaten.
Die \emph{Peer Interface FI}, die der direkten Verbindung mit weiteren Digitalen Zwillingen
dient, ist für einen Zwilling ohne Vernetzungsanforderung nicht relevant.

In ISO~23247 werden acht Typen beobachtbarer Fertigungselemente unterschieden,
sogenannte \emph{Observable Manufacturing Elements} (OME) \cite{iso23247}.
Fertigungsanlagen und Betriebsmittel werden dem Typ \emph{Equipment} zugeordnet,
die durch sie ausgeführten Operationen dem Typ \emph{Process}.
Das in der Anlage verarbeitete oder beförderte Werkstück ordnet die Norm dem Typ
\emph{Product} zu und trennt es damit konzeptuell von der Fertigungsanlage selbst.

\subsection{NVIDIA Omniverse-Ökosystem}
\label{subsec:omniverse}
NVIDIA Omniverse ist eine GPU-beschleunigte Kollaborationsplattform, die
Engineering-Werkzeuge verschiedener Hersteller über Konnektoren in OpenUSD-Szenen
zusammenführt, Änderungen über den Nucleus-Server im 24/7-Dauerbetrieb zwischen
allen Clients synchronisiert und Anlage sowie Produkt parallel in derselben Szene
simuliert \cite{nvidia_omniverse}.
Als Simulationsanwendung des Ökosystems verbindet \emph{Isaac Sim} physikalisch
korrekte Starrkörperdynamik mit fotorealistischem RTX-Rendering und wird in dieser
Arbeit für die Entwicklung der Digitalen Zwillinge eingesetzt \cite{nvidia_isaacsim}.
Auch in der Automobilindustrie wird Omniverse für die virtuelle Fabrikplanung
eingesetzt, um Planungsdaten verschiedener Werkzeuge ohne Kompatibilitätshürden in
einer gemeinsamen Simulationsumgebung zusammenzuführen \cite{bmw_nvidia_2021}.
Wie eine solche Simulation in Isaac Sim aussieht, zeigt Abbildung~\ref{fig:isaac_bmw}.

\begin{figure}[H]
    \centering
    \includegraphics[width=\linewidth]{BMW_Isaac Sim.jpg}
    \caption{NVIDIA Isaac Sim Viewport mit aktiviertem RTX Path Tracing. Simultane
        Darstellung unterschiedlicher Verarbeitungszustände in einer USD-Szene
        \cite{bmw_nvidia_2021}}
    \label{fig:isaac_bmw}
\end{figure}

Im Viewport sind gleichzeitig unterschiedliche Darstellungszustände einer USD-Szene
sichtbar.
Die Karosserie im Vordergrund ist vollständig mit RTX Path Tracing berechnet.
Reflexionen auf der Metalloberfläche und der Schattenwurf entstehen durch physikalisch
basierte Lichtpfadsimulation.
Der rechte Szenenbereich, bestehend aus Gabelstapler, Hallendecke und angrenzenden
Strukturen, wird im Wireframe-Modus dargestellt und repräsentiert approximierte
Kollisionsgeometrien.
Gabelstapler und Teile der Frontschürze sind mit einem weißen Standardmaterial
versehen, das Omniverse automatisch zuweist, wenn keine Materialdefinition im
Datensatz vorhanden ist.

\emph{Omniverse Kit} als Kern der Plattform ist ein erweiterbares
Anwendungsframework mit drei Eigenschaften, die es von klassischen
Fabrikplanungswerkzeugen unterscheiden.

\begin{itemize}
    \item \textbf{Konnektoreninfrastruktur:} Omniverse bindet bestehende
        Engineering-Werkzeuge über sogenannte \emph{Connectors} an, die
        herstellerspezifische Daten in OpenUSD übersetzen und einen bidirektionalen
        Datenaustausch herstellen. Dies umfasst Anbindungen an CAD-, BIM- und
        PLM-Systeme wie Autodesk Revit, Siemens NX und PTC Creo
        \cite{nvidia_omniverse_connectors}. Bestandsmodelle lassen sich so ohne
        manuellen Konvertierungsaufwand in einer gemeinsamen Szene zusammenführen.
    \item \textbf{Nucleus als zentrale Kollaborationsschicht:} Das technische Rückgrat
        bildet der \emph{Omniverse Nucleus} (\emph{database and collaboration engine}),
        an dem mehrere Clients gleichzeitig live verbunden sind. Änderungen an den
        Modellen werden in Echtzeit zwischen allen Clients synchronisiert, wobei die
        Plattform als Enterprise-Variante für einen hochverfügbaren 24/7-Dauerbetrieb
        in produktionsnahen Szenarien ausgelegt ist \cite{nvidia_nucleus_arch}.
    \item \textbf{Parallele Simulation von Anlage und Produkt:} Basierend auf dem
        NVIDIA DSX Blueprint für AI-Factory-Digitale-Zwillinge werden Geometrie-,
        thermische und elektrische Simulationsdaten verschiedener Disziplinen über
        USD-Layer verknüpft. Sowohl die Anlage als auch das prozessierte Produkt liegen
        in derselben Szene als parallel simulierbare Entitäten vor
        \cite{nvidia_omniverse}.
\end{itemize}

Isaac Sim wird mit einem vorinstallierten Extension-Katalog ausgeliefert, der
allgemeine Werkzeuge innerhalb der Simulationsumgebung bereitstellt, wie eine
VR-Experience oder den Script Editor, über den Python-Code direkt in der laufenden
Szene ausgeführt werden kann.

Das physikalische Fundament der Simulation bildet die Physik-Engine \emph{PhysX},
die Starrkörperdynamik, Gelenkkräfte und Kollisionserkennung GPU-parallel berechnet
\cite{nvidia_physx}.
PhysX ist zudem in verbreiteten Simulationsumgebungen wie Visual Components oder
KUKA.Sim integriert.

Zum Omniverse-Ökosystem gehören weitere spezialisierte Anwendungen.
\textbf{Isaac Lab} erweitert Isaac Sim um ein Framework für Reinforcement Learning,
mit dem Roboter-Verhaltensweisen trainiert und per Sim-to-Real-Transfer auf reale
Hardware übertragen werden können, wie in Abbildung~\ref{fig:isaac_lab} gezeigt
\cite{nvidia_isaaclab}.
Als drittes Glied der Isaac-Familie zeichnet \textbf{Isaac TeleOp} menschliche
Bedienabläufe als Trainingsdaten auf, wobei verschiedene Eingabegeräte wie etwa
VR-Brillen verwendet werden, wie in Abbildung~\ref{fig:isaac_teleop} gezeigt
\cite{nvidia_isaacteleoperation}.
\textbf{Isaac GR00T} stellt ein Grundlagenmodell für humanoide Roboter bereit und
kann aus den von Isaac TeleOp aufgezeichneten Bewegungsabläufen humanoide Roboter
trainieren, wie in den Abbildungen~\ref{fig:isaac_lab} und~\ref{fig:isaac_teleop}
gezeigt \cite{nvidia_groot}.
\textbf{Isaac ROS} ist für Mobilitätsanwendungen bis hin zum autonomen Fahren
ausgelegt \cite{nvidia_isaacros}.
Um eine ausreichende Datensatzvielfalt für KI-Modellanwendungen sicherzustellen,
wurde \textbf{NVIDIA Cosmos} entwickelt, um aus bestehenden Simulationsszenen
Variationen in Beleuchtung, Textur und Materialeigenschaften zu erzeugen
\cite{nvidia_cosmos}.

\begin{figure}[H]
    \centering
    \begin{minipage}[t]{0.48\linewidth}
        \centering
        \includegraphics[width=\linewidth]{isaac lab.png}
        \caption{Paralleles Training von Robotern in NVIDIA Isaac Lab
            \cite{nvidia_isaaclab}}
        \label{fig:isaac_lab}
    \end{minipage}
    \hfill
    \begin{minipage}[t]{0.48\linewidth}
        \centering
        \includegraphics[width=\linewidth]{NVIDIA Cosmos.png}
        \caption{Erfassen und Übertragen menschlicher Bedienabläufe per VR-Brille
            mit Isaac TeleOp und Isaac GR00T
            \cite{nvidia_isaacteleoperation,nvidia_groot}}
        \label{fig:isaac_teleop}
    \end{minipage}
\end{figure}


\subsection{Datenformate OpenUSD und URDF}
\label{subsec:datenformate}
Das von Pixar Animation Studios entwickelte Format \emph{Universal Scene Description}
(OpenUSD) organisiert Szenendaten als hierarchischen Namensraum aus \emph{Prims}
(Primitives) \cite{pixar_openusd}.
NVIDIA Omniverse nutzt OpenUSD als natives Dateiformat, sodass CAD-Geometrien,
Simulationsparameter und Laufzeitdaten verschiedener Planungsphasen ohne
werkzeugspezifische Konvertierungen in einer gemeinsamen Szene zusammenwachsen
\cite{nvidia_omniverse}.
Dies schafft auf Datenformatsebene die Grundlage für das durchgängige Datenmodell,
das VDI~4499 als Ziel der Digitalen Fabrik beschreibt.

Für die roboterspezifische Strukturbeschreibung ergänzt das \emph{Unified Robot
Description Format} (URDF) das OpenUSD-Ökosystem \cite{ros_urdf}.
URDF formalisiert kinematische Glieder (\emph{links}), Gelenke (\emph{joints}) sowie
visuelle, kollisions- und trägheitsrelevante Geometrien in einer XML-Struktur und wird
in der Robotik als verbreitetes Standardbeschreibungsformat eingesetzt \cite{ros_urdf}.


\subsection{Middleware-Architekturen für die industrielle Kommunikation}

Middleware-Architekturen bilden die kommunikative Brücke zwischen der physischen Anlage
und der Omniverse-Extension, die den digitalen Zwilling in der Simulationsszene steuert.
Sie übernehmen den Transport von Sensordaten, Steuerungsbefehlen und Zustandsinformationen
in Echtzeit.
Die im Folgenden vorgestellten Protokolle OPC~UA, MQTT und ROS~2 sind industriell
etabliert, weit verbreitet und decken unterschiedliche Anforderungen an Semantik,
Bandbreite und Echtzeitfähigkeit ab.

\subsubsection{OPC UA (DIN EN IEC 62541)}
Im Client-Server-Modell stellt die Steuerungseinheit einer Anlage als OPC~UA-Server
Maschinensignale und Zustandsdaten strukturiert bereit.
Ein Client liest Sensorwerte aus und schreibt bei Bedarf Sollwerte zurück.
Integrierte Sicherheitsmechanismen und plattformunabhängige Implementierung haben
OPC~UA zum verbreiteten Standard für die vertikale Integration in
Industrie~4.0-Architekturen gemacht.
Das Protokoll \emph{Open Platform Communications Unified Architecture} (OPC~UA) ist in
der Normenreihe DIN~EN~IEC~62541 festgeschrieben \cite{iec62541}.

\subsubsection{MQTT v5.0}
Im Publish-Subscribe-Ansatz veröffentlichen Sender Nachrichten an einen Broker,
der sie an alle Abonnenten des jeweiligen Themas weiterleitet, was MQTT~5.0
besonders für ressourcenbeschränkte IoT-Knoten geeignet macht.
Das Protokoll \emph{Message Queuing Telemetry Transport} (MQTT) ist in Version~5.0
von OASIS standardisiert \cite{oasis_mqtt5}.

\subsubsection{ROS 2 und DDS}
ROS\,2 ist für die Ansteuerung kinematischer Ketten ausgelegt, da Gelenkwinkel von
Roboterarmen über standardisierte Nachrichtentypen in Echtzeit übertragen werden.
Konfigurierbare \emph{Quality-of-Service}-Profile erlauben es, die Kommunikation gezielt
auf Latenz, Zuverlässigkeit oder Durchsatz abzustimmen, was die Kopplung von Isaac Sim
mit realer Roboterhardware unterstützt \cite{ros2_humble_docs}.
Praktische Einsatzbereiche reichen von der Intralogistik bis zur Fertigungsautomation
mit kameragestützten Pick-and-Place-Aufgaben.
Für Anwendungen im autonomen Fahren stellt NVIDIA mit Isaac ROS hardwarebeschleunigte
ROS\,2-Pakete bereit \cite{nvidia_isaacros}.
Die technische Grundlage bildet der \emph{Data Distribution Service} (DDS), auf dem
\emph{Robot Operating System~2} (ROS\,2) aufbaut \cite{ros2_humble_docs}.
